%% 
%% Template article for cas-sc documentclass for 
%% double column output.

\documentclass[a4paper,fleqn]{cas-sc}

% If the frontmatter runs over more than one page
% Use the longmktitle option.

%\documentclass[a4paper,fleqn,longmktitle]{cas-sc}

%\usepackage[numbers]{natbib}
%\usepackage[authoryear]{natbib}
\usepackage[numbers]{natbib}
%% The amssymb package provides various useful mathematical symbols
\usepackage{amssymb}
%% The amsmath package provides various useful equation environments.
\usepackage{amsmath}
\usepackage{epstopdf}
\usepackage{mathptmx,mathtools}
\usepackage{subcaption}
\usepackage{float} % Add this in your preamble
\usepackage[export]{adjustbox}
\usepackage{array} % Needed for m{} column type
\usepackage[percent]{overpic} % <-- put this in your preamble



\usepackage{amsmath}
\usepackage{nicefrac}
\usepackage{dirtytalk}
\usepackage{graphicx}   % For including images
\usepackage{tikz}
\usetikzlibrary{calc,tikzmark} % needed <<<<<<<<<<
\usepackage{xcolor}     % For coloring
\usepackage{lineno}
\renewcommand\linenumberfont{\normalfont\bfseries\small} % <-- Increase line number size
\usetikzlibrary{arrows.meta}
\usetikzlibrary{calc, positioning}
%%%Author macros
\def\tsc#1{\csdef{#1}{\textsc{\lowercase{#1}}\xspace}}
\tsc{WGM}
\tsc{QE}

\begin{document}

\setcounter{figure}{1}

	%%%%%%%%%%%%%%%%%%%%%%%%%%%%%________Figure_1___________%%%%%%%%%%%%%%%%%%%%%
\begin{figure}[!t]
	\centering
	% --------- Subfigure (b) --------- %
	\begin{overpic}[height=0.50\textheight,trim={1cm 0cm 1cm 0cm},clip]{Fig1a.eps}
		\put(5,97){\fontsize{13}{12}\selectfont\fontfamily{ptm}\selectfont\textbf{(a)}}
		
		\put(75,70){\fontfamily{ptm}\selectfont\textbf{Break-up}}
		\put(40,20){\fontfamily{ptm}\selectfont\textbf{Non-breaking}}
		
	\end{overpic}
	% --------- Subfigure (a) --------- %
	\begin{overpic}[height=0.50\textheight,trim={0cm 0cm 26cm 0cm},clip]{Fig1b.pdf}
		\put(5,101){\fontsize{13}{12}\selectfont\fontfamily{ptm}\selectfont\textbf{(b)}}
	\end{overpic}
	\caption{(a) Location of the simulated cases (C$_1$–C$_7$) from Table~\ref{Table1} on the Eötvös–Galilei (Eo–Ga) map for rising bubbles. The red dashed line indicates a constant Morton number (Mo = $10^{-3}$), while the blue dashed line divides break-up region from non-breaking region. (b) Schematic of the computational domain for rising bubble simulation on two-dimensional \textit{axisymmetric} with SI-LSM$_{\text{col}}$.}
	\label{Fig1}
\end{figure}
%%%%%%%%%%%%%%%%%%%%%%%%%%%%%________Figure_1_end___________%%%%%%%%%%%%%%%%%%%%%

\end{document}